The three subproblems have been connected through one coupled magnetic model. For P1, the proportional-damping Birkhoff representation, compatible endpoint projection and generating function lead to a dimensionally consistent weighted map. Ordinary midpoint remains a distinct implementation of the full vector field. Noncommuting matrix elimination retains the order of the mass, stiffness and damping factors, and the effective modal mass is derived from the stated normalization rather than substituted for the assembled beam mass.

For P2, integral observations identify the conservative residual and damping in a common coordinate system. The compensation retains the nominal cosine coupling and passes through the same physical force--torque and field--gradient interface in every control ablation. Its realized closed-loop equation keeps both the potential error and damping error. The restricted noiseless experiments test this complete connection; their small fitting residual is not a claim about general CHNN training or a calibrated physical device.

For P3, placing feedback inside the implicit midpoint equation gives convergence to the same actual closed loop under local smoothness, stability and domain assumptions. The error bound includes the accumulated increment residual, so an algebraic solver tolerance and an integration step are not interchangeable accuracy specifications. Fixed identification error produces a separate deviation from the nominal target. Exact symplecticity requires removal of the nonconservative residual. The weighted scheme for proportional residual damping preserves the corresponding conformal structure. The discrete KAM statement adds the necessary analytic, twist, nonresonance and smallness hypotheses; finite-time convergence and geometric diagnostics do not establish those hypotheses automatically.

The revised material and magnet record provides a consistent physical basis for these statements. Its endpoint inertia, modal scales, input conversion and energy-clearance bound are carried through the article. Extension to more retained modes must preserve the trained controller while accounting for the full tip angle and an explicitly defined full-model residual. An exact full-modal realization under stronger feedback assumptions is a separate theoretical construction. The current two-mode trajectories do not substitute for that staged validation or for an implemented energy governor.

The remaining empirical requirements are specimen identification, joint field--gradient and sensing calibration, and independent full-order or COMSOL tests. A finite-amplitude KAM certificate requires validated analytic remainders and theorem-specific constants on a declared domain. The article therefore supplies the scheme, controller, controlled convergence analysis and a reproducible verification sequence for Direction I, while keeping those further certification tasks explicit.
